\begin{statement}{E.44}{Proved}{The exponent-three class-three boundary quotient and its exact comparison with the class-two towers}
Retain the literal groups $H\triangleleft F\triangleleft B$,
the ten Schreier generators $\eta_i$, and the common source
$\mathscr D$ of E.24.  Put
\[
 \mathcal Q_3=B/\bigl(\gamma_4(H)H^3\bigr),\qquad
 \mathcal J_3=\operatorname{im}
       \bigl(B\longrightarrow\mathscr D\times\mathcal Q_3\bigr).
 \tag{E.44.1}
\]
The geometric kernel of $\mathcal Q_3\to B/H$ has order
$3^{175}$, with successive dimensions $10,45,120$.
Its actual projection to the class-two quotient $P_3$ of
E.33 forgets the last $120$ coordinates.

Every one of the $432$ multiplier-one source points of
E.24 has a generating lift to $\mathcal J_3$.
Over each such source point the complete finite fibre
of correction images has $3^{28}$ elements.
Exactly $3^6$ of its $3^{11}$ E.33 images lift to this
class-three quotient.  Each of those images has an
affine $\mathbb F_3^{22}$ fibre.

The comparison with the full three-adic class-two
tower of E.37 imposes four further independent
congruences on its eight free Smith coordinates.
They are explicitly \textup{(E.44.13)--(E.44.14)} below.
Over each fixed source point, $3^5$ of the $3^9$
E.37 reduction images satisfy them, and there are
$3^{27}$ compatible images in $\mathcal Q_3$.
The complete compatible inverse-limit fibre, retaining
the class-three coordinate, has the explicit continuous
parametrization
\[
                    \mathbb Z_3^{10}\times\mathbb F_3^{22}
 \tag{E.44.2}
\]
constructed in the proof.  It can be prescribed together
with every two-adic class-two family of E.41 and the
prime-to-six identity factors of E.42; one profinite
commutator word realizes each such family, and every
finite common source is generated.
The same families retain E.43's $\mathrm{SL}_2(\mathbb F_5)$
Ferrari refinement in the actual common sources
\textup{(E.44.15)} below.

There is a specified point of the entire E.37 tower,
over E.24 row $0$ with parameters $(0,0;0)$, that has
no lift to \textup{(E.44.1)}: take
\[
 u=w=0,\qquad Y_{37}=1,\qquad
                 Y_{38}=\cdots=Y_{44}=0
 \tag{E.44.3}
\]
in E.37's coordinates.  Thus the new obstruction
remains after all the earlier class-two levels have
been imposed.
\end{statement}
\begin{proof}
We first give the exact quotient, including its cube
relations.  Work in the truncated free associative
algebra over $\mathbb F_3$ on $X_1,\ldots,X_{10}$,
with monomials of degree at least four set to zero,
and send $\eta_i$ to $1+X_i$.
The group $M$ generated by these units is the quotient
\[
 H/\bigl(\gamma_4(H)\,\gamma_2(H)^3 H^9\bigr).
 \tag{E.44.4}
\]
Here and below a power subgroup means the subgroup
generated by the indicated powers, with closure in
the profinite notation.

To prove \textup{(E.44.4)}, ordinary class-three Hall
collection gives ten generator exponents modulo nine,
forty-five degree-two exponents modulo three, and
$330$ ordinary degree-three exponents modulo three.
Their collected forms are independent in the displayed
algebra.  The degree-one and degree-two terms first
recover the generator residues and commutator
exponents.  Changing a generator exponent by three
adds $X_i^3$ in degree three.  The remaining ordinary
degree-three Hall terms have the $330$ Lyndon leading
words, each with coefficient one, and are independent
of the ten $X_i^3$.  This proves both the inverse
collection map and the claimed kernel.
In particular $|M|=3^{395}$.

Every cube in $M$ is central.  For a unit $1+A$ its
cube is $1+v^3$, where $v$ is the degree-one part of
$A$.  The complete span of these cube values has
dimension $220$.  An explicit basis consists of
\[
 \begin{gathered}
 X_i^3,\qquad
 X_i^2X_j+X_iX_jX_i+X_jX_i^2\quad(i\ne j),\\
 \sum_{\sigma\in S_3}
       X_{\sigma(i)}X_{\sigma(j)}X_{\sigma(k)}
                           \quad(i<j<k).
 \end{gathered}
 \tag{E.44.5}
\]
Their disjoint monomial supports prove independence;
their numbers are $10,90,120$.
Each is an actual linear combination of cubes:
the repeated-index terms follow from
$(X_i+X_j)^3$ and $(X_i-X_j)^3$, and the
three-distinct-index term follows by subtracting
the three two-variable cubes from
$(X_i+X_j+X_k)^3$ and adding the single-variable cubes.
Conversely expansion of any $v^3$ belongs to this span.
Because the subgroup in \textup{(E.44.4)} lies in
$\gamma_4(H)H^3$, quotienting by precisely this central
span gives the required group
$H/(\gamma_4(H)H^3)$.

Here is a small coordinate model of that exact
quotient.  Retain E.33's coordinates
$v\in\mathbb F_3^{10}$ and
$z\in\Lambda^2\mathbb F_3^{10}$.
For $i<j<k$, let $t_{ijk}$ be the coefficient of
$X_iX_jX_k$ minus that of $X_jX_iX_k$.
This map annihilates \textup{(E.44.5)} and takes
$[X_i,[X_j,X_k]]$ to the corresponding unit vector.
On the $340$-dimensional restricted degree-three
term its kernel therefore has dimension $220$,
and is exactly the complete cube span just proved.
It gives a bijective group coordinate map
\[
 H/(\gamma_4(H)H^3) \longrightarrow\
 \mathbb F_3^{10}\oplus\mathbb F_3^{45}
                         \oplus\mathbb F_3^{120}.
 \tag{E.44.6}
\]

For explicit multiplication put, for $i<j$,
\[
 A(v,z)_{ij}=2v_iv_j+z_{ij},\qquad
 A(v,z)_{ji}=2v_iv_j-z_{ij},\qquad
 A(v,z)_{ii}=2(v_i^2-v_i).
 \tag{E.44.7}
\]
These are the actual quadratic Magnus coefficients.
For $(v,z,t)(w,a,b)=(v',z',t')$ the law is
\[
 \begin{split}
 v'&=v+w,\qquad z'=z+a+2v\wedge w,\\
 t'_{ijk}&=t_{ijk}+b_{ijk}
       +v_i A(w,a)_{jk}-v_j A(w,a)_{ik}\\
   &\hspace{24mm}
       +\bigl(A(v,z)_{ij}-A(v,z)_{ji}\bigr)w_k.
 \end{split}
 \tag{E.44.8}
\]
The inverse coordinate map is equally explicit.
With exponents represented by $0,1,2$, take the
ordered product of $\eta_i^{v_i}$, then of
$[\eta_i,\eta_j]^{z_{ij}-2v_iv_j}$, and finally
of $[\eta_i,[\eta_j,\eta_k]]^{c_{ijk}}$, where
\[
          c_{ijk}=t_{ijk}-v_iv_jv_k-v_i z_{jk}+v_j z_{ik}.
 \tag{E.44.9}
\]
The displayed product has exactly \textup{(E.44.6)}
as its coordinates.  This also specifies all
word preimages used by the computation.

The arithmetic source maps are the same literal
substitutions as in E.33:
$s:(x,y)\mapsto(y,x)$ and
$r:(x,y)\mapsto(y,(xy)^{-1})$.
Substitute the ten fixed Schreier words and collect
by \textup{(E.44.8)}.  Denote the resulting genuine
actions by $S$ and $R$.  They have orders two and
three and induce the old E.33 actions on $(v,z)$;
on the new centre they induce
$\Lambda^3 S_1$ and $\Lambda^3 R_1$.
The program retains each of the ten complete
generator images and verifies the actual words
$x=(r^{-1}s)^2$ and $y=(sr^{-1})^2$ against
conjugation by $x$ and $y$.

For $g\in H/(\gamma_4(H)H^3)$ the two equations are
literally $(Sg)g=1$ and
$g^{-1}R(g^{-1})R^2(g^{-1})=1$.
They define the map $s\mapsto sf$, $r\mapsto f^{-1}r$.
This map fixes $sr$, hence fixes $x$; it sends
$y=sxs$ to $f^{-1}yf$.
For $z=(xy)^{-1}=ryr^{-1}$ its image is
$hzh^{-1}$ with $h=f^{-1}R(f^{-1})$.
Thus these are the actual multiplier-one peripheral
maps, with their based conjugators.

The linear norm matrix on the new centre is
\[
 M_3=\begin{pmatrix}
 I+\Lambda^3S_1\\
 -I-\Lambda^3R_1-(\Lambda^3R_1)^2
 \end{pmatrix},\qquad
           \operatorname{rank}_{\mathbb F_3}M_3=98.
 \tag{E.44.10}
\]
Hence its kernel has dimension $22$.
For each of E.33's nine admissible abelian values,
write its central coordinate as $z_0+K_0 a$, where
$K_0$ is the saved $45$-by-$10$ basis of the
E.33 central norm kernel.
At fixed $v$ all remaining norm equations are
affine linear in $(a,t)$: this follows directly
from \textup{(E.44.8)} and the collected action.
The nine complete $240$-by-$130$ matrices have
rank $103$, and every affine system is solvable.
Thus each abelian value has $3^{27}$ solutions,
with five independent conditions on $a$ and an
affine $22$-dimensional triple fibre.
All nine matrices, their row operations, kernels
and particular solutions are reproduced by the
finite certificates below.  Each of the $432$
actual source cosets contains exactly three of
these abelian values by the literal E.33 source
comparison.  This proves the counts $3^{28}$ and
$3^6$ once the complete new source image is checked.

That image follows by actual words.  E.33 provides
forty-five adjusted normal relators $u_\ell$ in
$R_3\cap[H,H]$, with $R_3=\ker(B\to D_6)\cap F'$,
whose degree-two reductions form a basis.
For every pair $j<k$ choose the saved coefficients
so that $u_{jk}=\prod_\ell u_\ell^{a_\ell}$ has
degree-two coordinate $e_j\wedge e_k$ modulo three.
The words
\[
                 [\eta_i,u_{jk}^{4}],\qquad i<j<k,
 \tag{E.44.11}
\]
lie in $\ker(B\to\mathscr D)\cap\gamma_3(H)$
and map to the $120$ standard new central vectors.
Indeed fourth powers kill the exponent-two binary
image and retain the exponent-three target;
the ternary image is already trivial.
The uncomputed triple coordinate of $u_{jk}$
does not enter its commutator, since that layer
is central.  These words have zero class-two
image over every coefficient ring, because they
belong to $\gamma_3(H)$ integrally.
The exact relator indices, integral adjustments
and $45$ coefficients for every word in
\textup{(E.44.11)} are retained.
Together with the complete E.33 image, they prove
that the new common-source kernel image is its
full inverse image, with dimensions $6,45,120$.
There is no further source or $F'$ constraint
on the last two coordinates.

We next compare with the full E.37 tower.
Its two extra Smith conditions modulo three
and the five new conditions just computed have
joint rank six, in all nine abelian cases;
the joint affine systems are nonempty.
In the notation of E.37 put
\[
 c=c_f,\qquad \bar u=t_2\bmod3,\qquad
 b=(\bar Y_{37},\ldots,\bar Y_{44})^t\in\mathbb F_3^8.
 \tag{E.44.12}
\]
The complete additional criterion is
\[
                         A_c b=\beta(c,\bar u),
 \tag{E.44.13}
\]
where the three rank-four matrices are
\[
 \begin{gathered}
 A_0=\begin{pmatrix}
 1&0&0&0&0&0&0&0\\
 0&1&0&0&0&0&0&2\\
 0&0&1&0&0&0&2&0\\
 0&0&0&1&0&0&2&2
 \end{pmatrix},\qquad
 A_1=\begin{pmatrix}
 1&1&0&0&0&0&0&2\\
 0&0&1&0&0&0&2&0\\
 0&0&0&1&0&0&2&2\\
 0&0&0&0&0&1&2&2
 \end{pmatrix},\\[2mm]
 A_2=\begin{pmatrix}
 1&0&0&0&0&1&2&2\\
 0&1&0&0&0&2&1&0\\
 0&0&1&0&0&0&2&0\\
 0&0&0&1&0&0&2&2
 \end{pmatrix},\qquad
 \begin{array}{c|ccc}
 &\bar u=0&\bar u=1&\bar u=2\\ \hline
 c=0&(0,0,0,0)&(1,0,0,2)&(1,0,2,0)\\
 c=1&(2,0,2,0)&(1,2,0,0)&(2,0,0,2)\\
 c=2&(1,2,2,1)&(0,0,0,1)&(0,2,0,0)
 \end{array}
 \end{gathered}
 \tag{E.44.14}
\]
The table entries are the column vectors $\beta$.
These matrices are obtained by substituting the
actual E.37 origin and eight free Smith columns
into the class-three norm equations, and eliminating
only the new $120$ triple variables.
In particular they retain both old Smith factors
three; no arbitrary class-two solution has been
treated as a point of the inverse-limit tower.

For an explicit continuous parametrization choose
$u,w\in\mathbb Z_3$ and form the E.37 abelian point
$t_2=u$, $t_1=c-u+3w$, with its specified Hensel root
$t_3$.  In each matrix $A_c$ the four pivot columns
are unit columns.  Choose the four nonpivot Smith
coordinates freely in $\mathbb Z_3$; set the pivot
coordinates equal to the indicated affine solutions
of \textup{(E.44.13)}, using the digits $0,1,2$
for $\beta$, plus three times four further arbitrary
elements of $\mathbb Z_3$.
This is a continuous bijection from $\mathbb Z_3^{10}$
onto the compatible E.37 points, with inverse obtained
by recovering $u,w$, the four free coordinates and
dividing the four pivot discrepancies by three.
For each such point a saved particular solution of
\textup{(E.44.10)} and the $22$ saved kernel columns
give its complete affine triple fibre.
This proves \textup{(E.44.2)} with actual maps.
There are $3^4$ compatible old central reductions
at each of the three abelian values, hence $3^5$
per source and $3^{5+22}=3^{27}$ compatible
class-three images.

The comparison maps also prove simultaneity.
The common quotient of the new geometric group
and the class-two pro-three group is exactly
$H/(\gamma_3(H)H^3)$: every class-two quotient
of the former kills its new centre, and both
displayed maps onto this common quotient are
the literal projections.
More concretely, start with the E.37 word having
a prescribed class-two tower value.  Any desired
class-three value with the same $P_3$ projection
differs from its actual value by a new central
vector.  Correct this vector with the words
\textup{(E.44.11)}.  To retain all E.42 factors,
take their three-supported continuous powers.
This preserves their exponent-three images,
kills their images in every full pro-$p$ factor
for $p\ne3$, and still has exactly zero image in
every class-two quotient.  It also retains all
$\mathscr D$ coordinates.  This constructs one
profinite commutator word for every claimed family.

Generation is checked on the actual common source.
In $\mathcal J_3$ the image of $N_B$ is finite
nilpotent with primes two and three: its old image
has this property, and its new image is a
three-group, including the cyclic-three image in
$F/H$.  It is generated by the $97$ fixed
generators and maps onto $(\mathbb Z/6)^{97}$
through $G_6$.  Its Sylow Frattini dimensions are
therefore exactly $97$.  The retained $G_6$ map
is an automorphism on these quotients, proving
surjectivity on $N_B$ and then on $\mathcal J_3$.
The direct check of the induced $H$ abelian map,
using $x\mapsto x$, $y\mapsto f^{-1}yf$, has
rank ten for all nine abelian values as well.

For the simultaneous two-adic and three-adic
class-two factors, choose a lower-central index
larger than the nilpotence classes of the old
Ferrari $N_B$ image and every three-primary
image, including the new one.  Exactly as in
E.41, the image at that index supplies the pure
stable two-primary lower-central subgroup of
the boundary factor.  The corrected map is
already onto that factor, because its old value
has been retained.  Removing this pure subgroup
leaves the finite nilpotent common $N_B$ image,
so the same rank-$97$ Frattini argument applies.
For further prime-to-six factors the actual $H$
images have coprime orders to the existing
two-three image; their product over $B/H$ is
the actual common image by E.42's proved map.
The identity values on them give surjectivity.
No invariance of a fixed Ferrari source kernel
under the shadow is used.

For the simultaneous E.43 refinement, let $G_L$, $C$
and $S=G_L/C$ be its actual source, perfect kernel
and maximal solvable quotient.  Take a cofinal system
of the finite common boundary sources $G_{\rm old}$
just constructed that also contain $P_{20}$.
Their actual maps to $S$ are supplied by
$P_{20}\to S$ from E.43.  Their $H$ images are
solvable, including the new finite three-group;
the quotient $B/H$ is solvable as well.
Every common quotient with $G_L$ therefore factors
through $S$, and $S$ itself is common.  Thus the
literal source map gives precisely
\[
 \operatorname{im}\bigl(B\longrightarrow
          G_{\rm old}\times G_L\bigr)
       =G_{\rm old}\times_S G_L .
 \tag{E.44.15}
\]
For a word $f_0$ realizing a prescribed new boundary
family, its $S$ value is the E.43 value of the matching
Galois word $f_\sigma$, because all the old source
coordinates and $P_{20}$ have been retained.
The discrepancy
$\rho_L(f_0)^{-1}\rho_L(f_\sigma)$ lies in $C$.
The exact surjection
$\rho_L(R_H)=C$ of E.43, where
$R_H=\bigcap_{j\ge1}H^{[j]}$, supplies a single
$k\in R_H$ with that image.  Then $f=f_0k$
retains every old solvable $H$ image, including
the new class-three group, and agrees with
$f_\sigma$ in $G_L$.
The corrected image is onto both factors of
\textup{(E.44.15)} and has the same maximal common
quotient $S$, so it is the whole displayed fibre
product.  The four-strand Ferrari evaluation maps
are those of E.43 and are retained literally.
This proves the additional simultaneous assertion.

Finally the point \textup{(E.44.3)} has $t=0$
and doubled-log central vector $V_Me_{37}$ in
the exact Smith coordinates of E.37.
Its integral central norms vanish, so it is a
point of every class-two level.  It lies in the
specified actual source fibre, whose residue is
$c=0$.  The first row of $A_0$ at $\bar u=0$
requires $\bar Y_{37}=0$, contrary to
\textup{(E.44.3)}.
The certificate additionally records a row
annihilating all $120$ columns of $M_3$ and
evaluating to one on this point's norm error.
This proves the asserted failure of every
class-three lift over that tower point.

The three new programs in \texttt{code} have the shared
prefix \texttt{bc\_fermat\_class\_three\_} and suffixes
\texttt{burnside.py}, \texttt{norms.py}, and
\texttt{tower\_comparison.py}.  They
retain the complete quotient and action,
all nine affine systems, all $432$ source
records, the literal relator preimages,
the finite matrices in \textup{(E.44.14)}, and
the specified nonlifting point.  Their executed
outputs have status \texttt{PASS}; hashes verify
that every checkpoint05 input remained unchanged.
\end{proof}
